\section{Additional Details and Results}\label{sec:app_add_mod}
	\setcounter{table}{0}
	\renewcommand{\thetable}{H\arabic{table}}
	
	\setcounter{figure}{0}
	\renewcommand{\thefigure}{H\arabic{figure}}
\subsection{Microfoundation of perceived admission equations}\label{sec:microfound}
The subjective probability of regular admission conditional on taking the entrance exam is equal to the subjective probability that a student's believed score will be above the believed admission cutoff, over which the student forms a subjective probability distribution: $c_i^{Rb}\sim N(\bar{c}^{Rb},\sigma^2_{c^{Rb}})$. Letting $A^R_i$ denote a dummy for regular admission, the subjective probability of regular admission is:

\vspace{-12pt}
\begin{eqnarray}\label{eq:pr_Rb}
	Pr^b(A^R_i=1|\overline{PSU}^b_{i})&=&Pr\left( \overline{PSU}^b_{i}+\epsilon_i^{Pb}\geq \bar{c}^{Rb}+\epsilon_i^{c^{Rb}} \right)\\ \nonumber
	%&=&\Phi \left( \frac{\overline{PSU}^b_{i}-\bar{c}^{Rb}}{\sqrt{\sigma^{2}_{PSU^b}+\sigma^2_{c^{Rb}}}}\right) \\ \nonumber
	&=&\Phi\left(\gamma_0^b +\gamma_1^b \overline{PSU}^b_{i}\right),
\end{eqnarray}
\noindent where $\gamma_0^b=\frac{-\bar{c}^{Rb}}{\sqrt{\sigma^{2}_{PSU^b}+\sigma^2_{c^{Rb}}}}$ and $\gamma_1^b=\frac{1}{\sqrt{\sigma^{2}_{PSU^b}+\sigma^2_{c^{Rb}}}}$ and $\Phi(\cdot)$ is the standard Normal cumulative distribution function. Given an expected PSU score, uncertainty about admission is generated by uncertainty around own score ($\sigma^{2}_{PSU^b}$) and the admission cutoff ($\sigma^2_{c^{Rb}}$), which are absorbed by the parameters $\gamma_0^b$ and $\gamma_1^b$.\footnote{Several papers in the beliefs literature in Economics impose functional form restrictions on subjective probabilities (e.g. \citealp{delavande2019university,kapor2020heterogeneous}). We impose normality.} \par 


In treated schools, the subjective probability of preferential admission conditional on taking the entrance exam is equal to the subjective probability that a student's believed average GPA in the four high school years will be above the believed preferential admission cutoff, which is the $85^{th}$ percentile of high school GPA in the school. Students form a subjective probability distribution over the top 15\% cutoff in their school; we allow the mean of this distribution to vary across students: $c_i^{15b}\sim N(\bar{c}_i^{15b},\sigma^2_{c^{15b}})$. The subjective probability of preferential admission then is:
\begin{eqnarray}
	Pr^b(A^P_i=1|\overline{GPA}^{(9-12),b}_{i},\bar{c}_i^{15b})&=&Pr\left( \overline{GPA}^{(9-12),b}_{i}+\epsilon_i^{G^b}\geq c_0+\bar{c}_i^{15b}+\epsilon_i^{c15b} \right)\\ \nonumber
	&=&\Phi\left(\pi_0^b+\pi_1^b (\overline{GPA}^{(9-12),b}_{i}-\bar{c}_i^{15b})\right),
\end{eqnarray}
\noindent  where $\pi_0=\frac{-c_0}{\sqrt{\sigma^2_{GPA^b}+\sigma^2_{c15^{b}}}}$ and $\pi_1^b=\frac{1}{\sqrt{\sigma^2_{GPA^b}+\sigma^2_{c15^{b}}}}$.\footnote{Parameter $c_0$ is a net adjustment to the GPA and the cutoff to capture the fact that the top 15\% rule is based on an adjusted GPA measure, while the GPA and cutoff survey questions referred to unadjusted GPA to facilitate question comprehension.} Given an expected $GPA$ and an expected cutoff, uncertainty about admission is generated by uncertainty around own GPA ($\sigma^2_{GPA^b}$) and the school cutoff ($\sigma^2_{c15^{b}}$), which are absorbed by parameter $\pi_1^b$.



\subsection{College Entrants under Alternative Policies}\label{sec:app_add_mod2}
We analyze how the hypothetical interventions in Section \ref{sec:infointerventions} would affect the composition of college entrants and their persistence---an outcome of interest to  colleges. Figure \ref{fig:collegeentrantsinfo} describes college entrants under the no-intervention scenario (control condition) and the counterfactual interventions. Selection on test scores improves under the intervention that corrects all beliefs (third bar in Panel B). However, despite this improved selection, this intervention still results in college entrants with lower persistence rates (Panel A), because of substantial reductions in pre-college effort (Panel C). Unlike full belief correction, information on the importance of pre-college effort does not improve the baseline test scores of college entrants (fourth bar in Panel B), nor their pre-college effort (Panel C)---since the highest-achieving students (who are most likely to be admitted) tend to believe persistence is very likely, this information does not change much their beliefs---which helps explain why it increases pre-college effort in high school on average but does not substantially improve longer-term policy impacts. These findings show that when interventions affecting pre-college effort influence college persistence, they do so through two channels: changing the composition of students who enter college and altering the extent of their preparation during high school.\par

\begin{figure}[ht!]
    \centering
    \includegraphics[width=0.8\linewidth]{Revision/Figures R/college_entrants_info_count.png}
    \vspace{-0.5cm}
   % \captionsetup{skip=-0.4\baselineskip}
    \caption{Selective college persistence and characteristics of selective college entrants. This figure shows average 5-year college-persistence rates, baseline test scores, and pre-college effort for college entrants in the control group--where no intervention is introduced--and under three hypothetical interventions. The control group mean is represented by the first bar and by the dashed horizontal line. The three interventions are PACE (Baseline), and the two counterfactual interventions described in section \ref{sec:infointerventions}. SIMCE scores and pre-college effort are standardized to have mean zero and variance one in the control group. }
    \label{fig:collegeentrantsinfo}
\end{figure}